一般破缺规范理论的纵向等价关系~\eqnrefs{eq:goldstone-equivalence-general-dp11} 比较同一外部运动学下的纵向规范玻色子振幅与真空流形切向标量振幅。对 $W^+W^-$ 末态，定义精确差值
\begin{equation}
\boxed{
\mathscr R_{\rm GE}^{WW}(E,\theta)
\equiv
\mathcal M(W_L^-W_L^+;E,\theta)
-\mathcal M(\phi^-\phi^+;E,\theta).}
\label{eq:eeww-equivalence-dp11}
\end{equation}
在固定散射角、远离额外运动学奇点的高能缩放族上，树级等价定理给出
\begin{equation*}
 \lim_{m_Wc^2/E\to0}\mathscr R_{\rm GE}^{WW}(E,\theta)=0.
\end{equation*}
标准模型的 $\nu_e$、$\gamma$ 与 $Z$ 三类树图具有由规范对称性固定的相对系数；把它们在纵向极化高能极限中相加，表面的 $E^2/m_W^2c^4$ 增长严格相消。

取
\begin{equation*}
e^-(p_1)+e^+(p_2)\to W^-(k_-)+W^+(k_+),
\qquad s=(p_1+p_2)^2,
\end{equation*}
为单独检验高能手征相消，以下定义 $\me=0$ 的树级子模型。该模型中树级振幅必须同时包含 $t$ 道中微子交换和 $s$ 道光子、$Z$ 交换：
\begin{equation}
\mathcal M=\mathcal M_\nu+\mathcal M_\gamma+\mathcal M_Z.
\label{eq:eeww-three-amplitudes-dp11}
\end{equation}
三项使用同一套费米子电弱流和 $WWV$ Yang--Mills 顶角；它们的相对系数不是可独立调节的参数。

纵向极化在质心系中精确写为
\begin{equation*}
\epsilon_L^\mu(k)
=\frac1{m_Wc}\left(|\mathbf k|,\frac{E_W}{c}\widehat{\mathbf k}\right),
\qquad \epsilon_L\cdot k=0,
\end{equation*}
其中 $k^\mu=(E_W/c,|\mathbf k|\widehat{\mathbf k})$。定义差向量
\begin{equation}
\boxed{
 r_L^\mu(k)
 \equiv\epsilon_L^\mu(k)-\frac{k^\mu}{m_Wc}
 =\frac{m_Wc^2}{E_W+|\mathbf k|c}
 \left(-1,\widehat{\mathbf k}\right).}
\label{eq:longitudinal-polarization-highE-dp11}
\end{equation}
这给出精确分解 $\epsilon_L=k/(m_Wc)+r_L$，并有欧氏分量范数界
\begin{equation*}
 \|r_L(k)\|_2
 =\sqrt2\,\frac{m_Wc^2}{E_W+|\mathbf k|c}
 \le\sqrt2\,\frac{m_Wc^2}{E_W}.
\end{equation*}
因此可以先严格提取两个纵向外腿都取 $k/(m_Wc)$ 时的双动量系数，再把至少含一个 $r_L$ 的项保留在余项中。要证明规范相消，必须先证明不同树图的双动量系数落在同一个运动学结构上。定义无质量电子手征流与公共高能结构
\begin{equation}
\boxed{
J_\chi^\rho\equiv\bar v(p_2)\gamma^\rho P_\chi u(p_1),
\qquad
\mathcal K_\chi\equiv
\frac{J_\chi\!\cdot k_+}{m_W^2c^2},
\qquad \chi=L,R.}
\label{eq:eeww-common-kinematic-structure-dp75}
\end{equation}
由外腿 Dirac 方程有 $J_\chi\!\cdot(p_1+p_2)=0$，故也可用 $-J_\chi\!\cdot k_-$ 表示同一结构。Yang--Mills 三规范玻色子顶角的双动量收缩把两个 $s$ 道系数都写成某个耦合系数乘以 $\mathcal K_\chi$；左手中微子图的双动量收缩同样给出某个耦合系数乘以 $\mathcal K_L$。因此 $E^2/m_W^2c^4$ 系数的抵消是振幅级等式，而不只是两个耦合常数的数值巧合。

对右手电子，$W e\nu$ 顶角为零，所以中微子图消失。光子与 $Z$ 的主导项满足
\begin{equation}
\boxed{
\mathcal M_{\gamma,R}^{(E^2)}+\mathcal M_{Z,R}^{(E^2)}
=\left[-g^2s_W^2+g^2s_W^2\right]\mathcal K_R=0.}
\label{eq:eeww-right-cancellation-dp11}
\end{equation}
右手通道的高能抵消因而完全发生在 $\gamma$ 与 $Z$ 两个 $s$ 道振幅之间。

对左手电子，两个 $s$ 道图与 $t$ 道中微子交换分别给出
\begin{equation*}
\mathcal M_{L,s}^{(E^2)}=-\frac{g^2}{2}\mathcal K_L,
\qquad
\mathcal M_{L,t}^{(E^2)}=+\frac{g^2}{2}\mathcal K_L.
\end{equation*}
于是
\begin{equation}
\boxed{
\mathcal M_{L,s}^{(E^2)}+\mathcal M_{L,t}^{(E^2)}=0.}
\label{eq:eeww-left-total-cancellation-dp11}
\end{equation}
左右手通道采用不同图集却都消去表面的 $E^2/(m_W^2c^4)$ 增长，这同时检验了弱手征结构、弱混合角和三规范玻色子耦合之间的关系。

\begin{figure}[htbp]
\centering
\begin{tikzpicture}[line width=.8pt,line label/.style={fill=white,fill opacity=.96,text opacity=1,inner sep=1.15pt}]
\begin{scope}[xshift=-5.2cm]
\coordinate (a) at (-.8,0);\coordinate (b) at (.8,0);
\draw[fermion] (-2.4,1)--(a) node[line label,midway,above left]{$e^-$};
\draw[fermion] (a)--(-2.4,-1) node[line label,midway,below left]{$e^+$};
\draw[fermion] (a)--(b) node[line label,midway,above]{$\nu_e$};
\draw[decorate,decoration={snake,amplitude=1.2pt,segment length=5pt}] (b)--(2.4,1) node[right]{$W^-$};
\draw[decorate,decoration={snake,amplitude=1.2pt,segment length=5pt}] (b)--(2.4,-1) node[right]{$W^+$};
\node at (0,-1.7){$t$ channel};
\end{scope}
\begin{scope}[xshift=0cm]
\coordinate (a) at (-.8,0);\coordinate (b) at (.8,0);
\draw[fermion] (-2.4,1)--(a);\draw[fermion] (a)--(-2.4,-1);
\draw[decorate,decoration={snake,amplitude=1.2pt,segment length=5pt}] (a)--(b) node[line label,midway,above]{$\gamma$};
\draw[decorate,decoration={snake,amplitude=1.2pt,segment length=5pt}] (b)--(2.4,1);
\draw[decorate,decoration={snake,amplitude=1.2pt,segment length=5pt}] (b)--(2.4,-1);
\node at (0,-1.7){$s$ channel};
\end{scope}
\begin{scope}[xshift=5.2cm]
\coordinate (a) at (-.8,0);\coordinate (b) at (.8,0);
\draw[fermion] (-2.4,1)--(a);\draw[fermion] (a)--(-2.4,-1);
\draw[decorate,decoration={snake,amplitude=1.2pt,segment length=5pt}] (a)--(b) node[line label,midway,above]{$Z$};
\draw[decorate,decoration={snake,amplitude=1.2pt,segment length=5pt}] (b)--(2.4,1);
\draw[decorate,decoration={snake,amplitude=1.2pt,segment length=5pt}] (b)--(2.4,-1);
\node at (0,-1.7){$s$ channel};
\end{scope}
\end{tikzpicture}
\caption{$e^+e^-\to W^+W^-$ 的三类树图。删去规范对称性要求的任意一类贡献，纵向 $W$ 振幅都会留下非物理的高能幂次增长。}
\label{fig:eeww-gauge-cancellation-dp11}
\end{figure}

\begin{derivation}[从\eqnrefs{eq:eeww-three-amplitudes-dp11}到\eqnrefs{eq:eeww-right-cancellation-dp11}]
令中性传播子的入射四动量为 $q=p_1+p_2=k_-+k_+$。对三个动量都取向顶角流入的 Yang--Mills 张量，散射图中需要
\begin{equation*}
\Gamma_{\rho\mu\nu}(-q,k_-,k_+)
=(k_--k_+)_\rho g_{\mu\nu}
 +(k_++q)_\mu g_{\nu\rho}
 -(q+k_-)_\nu g_{\rho\mu}.
\end{equation*}
两个末态 $W$ 在壳，$k_-^2=k_+^2=m_W^2c^2$，并且
\begin{equation*}
q^2=s=2m_W^2c^2+2k_-\!\cdot k_+.
\end{equation*}
由精确分解~\eqnrefs{eq:longitudinal-polarization-highE-dp11}，完整双纵向收缩恒等地分成 $kk$、$kr_L$、$r_Lk$ 与 $r_Lr_L$ 四部分。先提取 $kk$ 部分，即在两个外腿上分别取 $k_-^\mu/(m_Wc)$ 与 $k_+^\nu/(m_Wc)$；顶角的三个张量项分别给
\begin{align*}
\Gamma_{\rho\mu\nu}k_-^\mu k_+^\nu
={}&(k_-\!\cdot k_+)(k_--k_+)_\rho\\
&+\bigl[k_-\!\cdot(k_++q)\bigr]k_{+\rho}
-\bigl[k_+\!\cdot(q+k_-)\bigr]k_{-\rho}.
\end{align*}
利用 $q=k_-+k_+$，两方括号都是 $m_W^2c^2+2k_-\!\cdot k_+$，于是
\begin{align*}
\Gamma_{\rho\mu\nu}k_-^\mu k_+^\nu
&=\bigl(m_W^2c^2+k_-\!\cdot k_+\bigr)(k_+-k_-)_\rho\\
&=\frac{s}{2}(k_+-k_-)_\rho .
\end{align*}
无质量电子的手征流满足
\begin{equation*}
q_\rho J_\chi^\rho
=\bar v(p_2)(\slashed p_1+\slashed p_2)P_\chi u(p_1)=0,
\end{equation*}
其中分别用了 $\slashed p_1u(p_1)=0$ 与 $\bar v(p_2)\slashed p_2=0$。因此 $J_\chi\!\cdot k_-=-J_\chi\!\cdot k_+$，从而
\begin{equation*}
\frac{1}{s}J_\chi^\rho
\Gamma_{\rho\mu\nu}k_-^\mu k_+^\nu
=J_\chi\!\cdot k_+.
\end{equation*}
光子图与 $Z$ 图的双 $k/(m_Wc)$ 项具有同一个运动学因子 $\mathcal K_\chi$。所有至少含一个 $r_L$ 的项由~\eqnrefs{eq:longitudinal-polarization-highE-dp11} 的范数界控制，在固定角高能极限中消失。

对右手电子，$P_Lu_R=0$，所以 $t$ 道中微子图严格消失。由
\begin{equation*}
g_{WW\gamma}=gs_W=g_{\rm em},\qquad
 g_{WWZ}=gc_W,\qquad
 g_{\gamma e_R}=-g_{\rm em},\qquad
 g_{Ze_R}=\frac{g}{c_W}s_W^2,
\end{equation*}
两个 $s$ 道主导项之和为
\begin{align*}
\mathcal M_{\gamma,R}^{(E^2)}+\mathcal M_{Z,R}^{(E^2)}
&=\left[g_{WW\gamma}g_{\gamma e_R}+g_{WWZ}g_{Ze_R}\right]\mathcal K_R\\
&=\left[-g^2s_W^2+g^2s_W^2\right]\mathcal K_R=0,
\end{align*}
即正文\eqnrefs{eq:eeww-right-cancellation-dp11}。这里的相消同时使用了 $e=gs_W$、$WWZ$ 的非阿贝尔顶角系数和右手电子的弱同位旋 $T^3=0$；任意改变其中一项都会留下随能量增长的非物理振幅。
\end{derivation}

\begin{derivation}[从\eqnrefs{eq:eeww-three-amplitudes-dp11}到\eqnrefs{eq:eeww-left-total-cancellation-dp11}]

左手中性流耦合为
\begin{equation*}
g_{Ze_L}=\frac g{c_W}\left(-\frac12+s_W^2\right).
\end{equation*}
由正文\eqnrefs{eq:eeww-common-kinematic-structure-dp75} 前的 Yang--Mills 双纵向收缩，两个 $s$ 道图共享 $\mathcal K_L$，故
\begin{align*}
\mathcal M_{L,s}^{(E^2)}
&=\left[g_{WW\gamma}(-g_{\rm em})+g_{WWZ}g_{Ze_L}\right]\mathcal K_L\\
&=\left[-g^2s_W^2+g^2\left(-\frac12+s_W^2\right)\right]\mathcal K_L\\
&=-\frac{g^2}{2}\mathcal K_L.
\end{align*}

再看中微子图。取 $r=p_1-k_-$，其张量部分为
\begin{equation*}
\mathcal M_\nu^{\mu\nu}
=-\frac{g^2}{2}\bar v(p_2)\gamma^\nu P_L
\frac{\slashed r}{r^2}\gamma^\mu P_Lu(p_1).
\end{equation*}
同样提取中微子图的双 $k/(m_Wc)$ 系数，分别取 $k_-^\mu/(m_Wc)$ 与 $k_+^\nu/(m_Wc)$。由 $\slashed p_1u(p_1)=0$，
\begin{equation*}
\slashed k_-P_Lu=(\slashed p_1-\slashed r)P_Lu
=-\slashed rP_Lu.
\end{equation*}
因此传播子分子与右侧收缩给出的 $\slashed r$ 合并，利用 $\slashed r\slashed r=r^2$ 消去传播子分母：
\begin{align*}
\mathcal M_{L,t}^{(E^2)}
&=\frac{g^2}{2m_W^2c^2}
\bar v(p_2)\slashed k_+P_Lu(p_1)\\
&=\frac{g^2}{2}\mathcal K_L.
\end{align*}
$\mathcal K_L$ 正是两个 $s$ 道已经出现的同一电子流结构。于是双动量系数满足
\begin{equation*}
\mathcal M_{L,s}^{(E^2)}+\mathcal M_{L,t}^{(E^2)}
=\left(-\frac{g^2}{2}+\frac{g^2}{2}\right)\mathcal K_L=0,
\end{equation*}
得到正文\eqnrefs{eq:eeww-left-total-cancellation-dp11}。
\end{derivation}
